\documentclass[11pt]{article}
\usepackage[margin=1in]{geometry}
\usepackage[T1]{fontenc}
\usepackage[utf8]{inputenc}
\usepackage{hyperref}
\usepackage{longtable}
\usepackage{booktabs}
\usepackage{xcolor}
\usepackage{array}
\usepackage{enumitem}
\usepackage{fancyvrb}
\usepackage{listings}

\lstset{
  basicstyle=\ttfamily\small,
  breaklines=true,
  frame=single,
  backgroundcolor=\color{gray!8},
  xleftmargin=4pt,
  xrightmargin=4pt
}

\setlength{\parindent}{0pt}
\setlength{\parskip}{0.7em}

\definecolor{passgreen}{RGB}{0,128,0}
\definecolor{failred}{RGB}{180,0,0}
\definecolor{warnamber}{RGB}{180,100,0}
\definecolor{infoBlue}{RGB}{0,70,160}

\newcommand{\pass}{\textcolor{passgreen}{\textbf{PASS}}}
\newcommand{\fix}{\textcolor{failred}{\textbf{FIX REQUIRED}}}
\newcommand{\improve}{\textcolor{warnamber}{\textbf{IMPROVE}}}
\newcommand{\new}{\textcolor{infoBlue}{\textbf{NEW}}}
\newcommand{\rfile}[1]{\texttt{#1}}
\newcommand{\rcode}[1]{\texttt{#1}}

\title{Version 2.0 Execution Plan\\Long Island Clustering Workflow\\
\large Phase-by-Phase Code Correction and Improvement Roadmap\\
Derived from \texttt{Li\_Clustering\_Audit\_Report.tex}}
\author{AMS 597 Project Team}
\date{\today}

\begin{document}
\maketitle
\tableofcontents
\newpage

% ============================================================
\section{Purpose and Scope of Version 2.0}
% ============================================================

This document is the authoritative implementation plan for producing
\rfile{li\_clustering\_v2.Rmd} from the existing \rfile{li\_clustering.Rmd}.
It translates every defect and improvement identified in
\rfile{Li\_Clustering\_Audit\_Report.tex} into precise, line-level code changes
organized by phase.

\subsection{What Version 2.0 Is and Is Not}

Version 2.0 is a corrected and strengthened copy of the existing workflow.
It is not a rewrite. Every phase that passed the audit is carried over
unchanged. Only the components identified as defective or improvable are
modified. The clustering methodology, feature engineering approach, and overall
pipeline structure remain the same.

\subsection{Pre-work: Create the v2 File}

Before making any changes, copy the source file:

\begin{lstlisting}
# In terminal from the project root:
cp topic3_nyc_clustering/li_clustering.Rmd \
   topic3_nyc_clustering/li_clustering_v2.Rmd
\end{lstlisting}

All changes described in this document are applied to
\rfile{li\_clustering\_v2.Rmd} only. The original
\rfile{li\_clustering.Rmd} is left untouched as the v1 reference.

\subsection{Change Inventory Summary}

\begin{longtable}{|p{1.2cm}|p{3.5cm}|p{2cm}|p{6.5cm}|}
\hline
\textbf{Phase} & \textbf{Change} & \textbf{Priority} & \textbf{What changes} \\
\hline
\endfirsthead
\hline
\textbf{Phase} & \textbf{Change} & \textbf{Priority} & \textbf{What changes} \\
\hline
\endhead
1 & Profile upgrade & Critical & \rcode{run\_profile <- "consensus\_16gb"} \\
\hline
1 & Version header & Low & Add v2.0 label to Rmd YAML title \\
\hline
3 & Binary flag encoding & Low & Replace \rcode{factor(..., labels = c("FALSE","TRUE"))} with \rcode{as.factor()} idiom \\
\hline
3 & Thunderstorm grouping & Low & Move Thunderstorm from grepl to explicit lookup \\
\hline
4 & Feature set sensitivity & High & Add reduced feature matrix (drop \rcode{Pressure(in)} and \rcode{Visibility(mi)}) \\
\hline
5 & Elbow plot annotation & Low & Add $R^2$ (variance explained) column to fit summary \\
\hline
6 & k-selection diagnostics & High & Add verbose logging block before and after fallback condition \\
\hline
6 & Quality gate warning & Critical & Add explicit \rcode{warning()} + printed block when Jaccard or silhouette below threshold \\
\hline
6 & Jaccard CI & Medium & Under consensus profile, compute SD of Jaccard across repeats \\
\hline
8 & Cluster 2 naming & Medium & Extend naming hierarchy to cover generic midday surface-road pattern \\
\hline
8 & Severity 4 explicit column & Medium & Add \rcode{sev4\_rate} to severity profile table with baseline row \\
\hline
9 & Stability plot annotation & Medium & Add subtitle to stability chart disclosing red-zone result \\
\hline
9 & Feature sensitivity comparison table & High & Export silhouette / Jaccard for full vs.\ reduced feature set \\
\hline
10 & Limitations table & Critical & Add explicit row naming cluster quality failure and exploratory status \\
\hline
7 & DBSCAN eps sensitivity & Medium & Run two additional trials at eps $= 0.9$ and eps $= 1.1$ \\
\hline
7 & Corridor centroid print & Medium & Print centroid lat/lng of each spatial cluster for manual verification \\
\hline
\end{longtable}

% ============================================================
\section{Phase 1: Profile and YAML Changes}
% ============================================================

\subsection{Change 1-A: Upgrade Run Profile (Critical)}

\textbf{File:} \rfile{li\_clustering\_v2.Rmd}\\
\textbf{Location:} Line 135 (Phase 1 config chunk \rcode{phase-1-config})\\
\textbf{Audit source:} Defect 3 in \rfile{Li\_Clustering\_Audit\_Report.tex}

\textbf{Problem:} \rcode{run\_profile <- "local\_safe"} uses a 5{,}000-row
sample (15\% of data) for validation with one bootstrap repeat.
Silhouette and Jaccard estimates from 15\% of the data are unreliable.
The observed Jaccard of 0.1989 may be a sampling artifact.

\textbf{Old code (line 135):}
\begin{lstlisting}
run_profile <- "local_safe"
\end{lstlisting}

\textbf{New code:}
\begin{lstlisting}
run_profile <- "consensus_16gb"
\end{lstlisting}

\textbf{Effect under consensus\_16gb:}
\begin{itemize}
\item Validation sample: 4{,}000 rows, 4 repeats, result averaged.
\item Stability sample: 2{,}500 rows, 3 repeats, 10 bootstraps.
\item Silhouette and Jaccard will be consensus estimates rather than
  single-sample point estimates.
\end{itemize}

\textbf{Memory note:} Under \rcode{consensus\_16gb}, peak memory during
Phase 6 should not exceed $\approx$8--10 GB. If the machine has less than
12 GB available to R, set \rcode{mem.maxVSize("12G")} at the top of the
setup chunk before the Phase 6 chunk runs.

\subsection{Change 1-B: Version Label in YAML (Low)}

\textbf{Location:} Lines 1--6, YAML header.

\textbf{Old code:}
\begin{lstlisting}
---
title: "Long Island Clustering Workflow"
author: "AMS 597 Project Team"
date: "`r format(Sys.Date(), '%B %d, %Y')`"
output: pdf_document
---
\end{lstlisting}

\textbf{New code:}
\begin{lstlisting}
---
title: "Long Island Clustering Workflow v2.0"
subtitle: "Corrected and Strengthened --- Based on Audit Report"
author: "AMS 597 Project Team"
date: "`r format(Sys.Date(), '%B %d, %Y')`"
output:
  pdf_document:
    toc: true
    number_sections: true
---
\end{lstlisting}

% ============================================================
\section{Phase 3: Feature Engineering Corrections}
% ============================================================

\subsection{Change 3-A: Binary Flag Factor Encoding (Low)}

\textbf{Location:} Lines 584--590 (chunk \rcode{phase-3-binary-flags})\\
\textbf{Audit source:} Improvement IP-7

\textbf{Problem:} \rcode{factor(as.logical(x), levels = c(FALSE, TRUE), labels = c("FALSE", "TRUE"))}
is unconventional. The string labels \rcode{"FALSE"} and \rcode{"TRUE"} are
indistinguishable by eye from logical values but are character strings,
which can cause confusion in \rcode{kproto()} distance computation and in
downstream filtering.

\textbf{Old code (lines 584--590):}
\begin{lstlisting}
for (flag_name in retained_flags) {
  df_phase3[, (flag_name) := factor(
    as.logical(get(flag_name)),
    levels = c(FALSE, TRUE),
    labels = c("FALSE", "TRUE")
  )]
}
\end{lstlisting}

\textbf{New code:}
\begin{lstlisting}
for (flag_name in retained_flags) {
  df_phase3[, (flag_name) := factor(
    as.logical(get(flag_name)),
    levels = c(FALSE, TRUE)
  )]
}
\end{lstlisting}

\textbf{Rationale:} Omitting the \rcode{labels} argument causes R to use
the native \rcode{FALSE} and \rcode{TRUE} factor level names, which are
standard and unambiguous. The \rcode{kproto()} function handles factor
columns by level count; the label string values do not affect distance
computation, but removing the string-label confusion improves code
transparency and avoids surprising \rcode{=="FALSE"} comparisons elsewhere.

\textbf{Propagation:} The same pattern appears at lines 628--636 in
Phase 4 where \rcode{weekend} and \rcode{rush\_hour} are converted to
factor. Apply the identical fix there:

\textbf{Old code (lines 627--636):}
\begin{lstlisting}
df_scenario_base[, weekend := factor(
  weekend,
  levels = c(FALSE, TRUE),
  labels = c("FALSE", "TRUE")
)]
df_scenario_base[, rush_hour := factor(
  rush_hour,
  levels = c(FALSE, TRUE),
  labels = c("FALSE", "TRUE")
)]
\end{lstlisting}

\textbf{New code:}
\begin{lstlisting}
df_scenario_base[, weekend := factor(weekend, levels = c(FALSE, TRUE))]
df_scenario_base[, rush_hour := factor(rush_hour, levels = c(FALSE, TRUE))]
\end{lstlisting}

\subsection{Change 3-B: Thunderstorm Weather Group (Low)}

\textbf{Location:} \rcode{weather\_group\_map} list in the helpers chunk (lines 18--33)
and the weather-grouping derivation in Phase 3 (chunk \rcode{phase-3-weather}).\\
\textbf{Audit source:} Improvement IP-8

\textbf{Problem:} Thunderstorm variants are matched by \rcode{grepl()} as a
fallback rather than being listed explicitly in \rcode{weather\_group\_map}.
Variants with spacing differences (e.g., ``Thunder Storm'', ``Thunderstorm / Windy'')
could be misclassified.

\textbf{New code --- update \rcode{weather\_group\_map} in the helpers chunk:}
\begin{lstlisting}
weather_group_map <- list(
  Clear      = c("Fair", "Clear", "Fair / Windy"),
  Cloudy     = c(
    "Mostly Cloudy", "Partly Cloudy", "Scattered Clouds", "Overcast",
    "Cloudy", "Mostly Cloudy / Windy", "Partly Cloudy / Windy",
    "Cloudy / Windy"
  ),
  Rain       = c(
    "Light Rain", "Rain", "Heavy Rain", "Light Drizzle",
    "Light Rain / Windy", "Heavy Rain / Windy", "Drizzle"
  ),
  Snow_Ice   = c(
    "Light Snow", "Snow", "Heavy Snow", "Wintry Mix",
    "Light Freezing Rain", "Ice Pellets", "Sleet"
  ),
  Fog_Haze   = c(
    "Fog", "Haze", "Smoke", "Mist", "Patches of Fog", "Shallow Fog"
  ),
  Thunderstorm = c(
    "Thunderstorm", "Thunder", "Thunder / Windy",
    "Thunderstorm / Windy", "Heavy Thunderstorm",
    "Light Thunderstorm", "Thunderstorms and Rain"
  )
)
\end{lstlisting}

The \rcode{grepl()} Thunderstorm pattern match in the Phase 3 chunk can be
retained as a final safety net \textit{after} the explicit lookup runs, but
should only catch truly novel strings rather than serving as the primary
classification path. Add a comment making this clear.

% ============================================================
\section{Phase 4: Reduced Feature Set Matrix (New)}
% ============================================================

\subsection{Change 4-A: Add Reduced Feature Scenario Matrix (High)}

\textbf{Location:} After the existing Phase 4 matrix construction
(after line $\approx$760, before Phase 5 begins).\\
\textbf{Audit source:} Improvement IP-1

\textbf{Problem:} \rcode{Pressure(in)} and \rcode{Visibility(mi)} both required
median imputation. Including imputed values in a distance-based clustering
method can pull imputed points artificially toward the global median,
deflating inter-cluster separation. The v2 workflow tests whether removing
these two features improves silhouette and Jaccard.

\textbf{New chunk to insert --- name it \rcode{phase-4-reduced-scenario-matrix}:}
\begin{lstlisting}
# v2.0: Reduced feature matrix for sensitivity analysis.
# Drops Pressure(in) and Visibility(mi) which required median imputation.
reduced_numeric_cols <- setdiff(
  scenario_numeric_cols,
  c("Pressure(in)", "Visibility(mi)")
)

df_scenario_reduced <- df_scenario[
  , c(scenario_factor_cols, reduced_numeric_cols), with = FALSE
]

# Verify: no NAs, correct column count
stopifnot(sum(is.na(df_scenario_reduced)) == 0)
stopifnot(ncol(df_scenario_reduced) ==
            length(scenario_factor_cols) + length(reduced_numeric_cols))

reduced_matrix_summary <- data.table(
  item  = c("Rows", "Factor columns", "Numeric columns",
            "Dropped from full matrix"),
  value = c(
    nrow(df_scenario_reduced),
    length(scenario_factor_cols),
    length(reduced_numeric_cols),
    paste(c("Pressure(in)", "Visibility(mi)"), collapse = ", ")
  )
)
knitr::kable(reduced_matrix_summary,
             caption = "Phase 4 reduced scenario matrix summary")
\end{lstlisting}

This matrix (\rcode{df\_scenario\_reduced}) will be used in the new Phase 5
reduced-feature candidate runs. Keep \rcode{df\_scenario} (full feature set)
intact for the primary pipeline.

% ============================================================
\section{Phase 5: k-Prototypes Additions}
% ============================================================

\subsection{Change 5-A: Add Variance-Explained Column to Fit Summary (Low)}

\textbf{Location:} Phase 5 fit diagnostics chunk (lines $\approx$858--918).\\
\textbf{Audit source:} Phase 5 issues noted in Audit Report.

\textbf{Problem:} The fit summary reports \rcode{tot.withinss} but not
variance explained ($R^2$). Without $R^2$ the elbow plot has no reference
scale; a reader cannot tell whether $k=5$ explains 40\% or 70\% of variance.

\textbf{Change:} In the block that builds \rcode{kproto\_fit\_summary},
add a \rcode{variance\_explained\_pct} column:

\begin{lstlisting}
# After computing kproto_fit_summary, add R^2 column:
total_ss <- sum(
  sapply(scenario_numeric_cols, function(col) {
    x <- df_scenario[[col]]
    sum((x - mean(x, na.rm = TRUE))^2, na.rm = TRUE)
  })
)

kproto_fit_summary[, variance_explained_pct := round(
  100 * (1 - tot_withinss / total_ss), 2
)]
\end{lstlisting}

\textbf{Note:} This is an approximation because \rcode{tot.withinss} in
k-prototypes includes both numeric and categorical components. The $R^2$
computed here applies only to the numeric component and should be
labelled as ``approximate numeric $R^2$'' in the report.

\subsection{Change 5-B: Add Reduced-Feature Candidate Runs (High)}

\textbf{Location:} After the existing Phase 5 candidate runs block.

\textbf{New chunk --- name it \rcode{phase-5-reduced-kproto}:}
\begin{lstlisting}
# v2.0: Run k-prototypes on the reduced feature matrix.
# Compares cluster quality when imputed features are excluded.

kp_results_reduced <- list()
for (k_val in candidate_k_values) {
  set.seed(42)
  kp_results_reduced[[as.character(k_val)]] <- clustMixType::kproto(
    x      = df_scenario_reduced,
    k      = k_val,
    lambda = NULL,
    nstart = 10,
    type   = "gower",
    keep.data = TRUE,
    verbose = FALSE
  )
}

reduced_fit_summary <- data.table(
  k = candidate_k_values,
  tot_withinss = sapply(
    candidate_k_values,
    function(k) kp_results_reduced[[as.character(k)]]$tot.withinss
  ),
  feature_set = "reduced"
)

full_fit_summary_labeled <- copy(kproto_fit_summary[, .(k, tot_withinss)])
full_fit_summary_labeled[, feature_set := "full"]

feature_comparison_fits <- rbind(
  full_fit_summary_labeled, reduced_fit_summary
)

fwrite(feature_comparison_fits,
       file.path(output_dir, "phase5_feature_comparison_fits.csv"))

knitr::kable(feature_comparison_fits,
             caption = "Phase 5 full vs reduced feature set within-cluster distances")
\end{lstlisting}

% ============================================================
\section{Phase 6: Validation and Stability Corrections}
% ============================================================

\subsection{Change 6-A: k-Selection Diagnostic Logging (High)}

\textbf{Location:} Lines 1284--1329 (chunk \rcode{phase-6-select-k}).\\
\textbf{Audit source:} Defect 2 in \rfile{Li\_Clustering\_Audit\_Report.tex}

\textbf{Problem:} The current fallback condition silently selects $k$ with
no printed trace of what was evaluated, what conditions triggered, and
why the final $k$ was chosen. When the fallback condition involves
\rcode{NA} extraction or data.table key mismatch, the code fails silently
and a different $k$ is selected without any diagnostic message.

\textbf{Old code (lines 1284--1289):}
\begin{lstlisting}
selected_k <- validation_table$k[1]
selected_stability <- stability_summary[k == selected_k, jaccard]

if (length(selected_stability) == 1 &&
    selected_stability < 0.5 &&
    nrow(stability_summary) > 1) {
  selected_k <- stability_summary[order(-jaccard, -rand, k)][1, k]
}
\end{lstlisting}

\textbf{New code (replace lines 1284--1289):}
\begin{lstlisting}
# --- v2.0: Verbose k-selection with explicit diagnostics ---
validation_optimal_k_raw  <- validation_table$k[1]
silhouette_at_optimal     <- validation_table$silhouette[1]

cat("=== Phase 6 k-Selection Diagnostics ===\n")
cat("Silhouette-optimal k        :", validation_optimal_k_raw, "\n")
cat("Silhouette at optimal k     :", round(silhouette_at_optimal, 4), "\n")
cat("Silhouette threshold (0.30) :",
    ifelse(silhouette_at_optimal >= 0.30, "MET", "NOT MET"), "\n\n")
cat("Stability summary available for k values:",
    paste(stability_summary$k, collapse = ", "), "\n")

selected_stability <- stability_summary[
  k == validation_optimal_k_raw, jaccard
]
cat("Jaccard at silhouette-optimal k :",
    ifelse(length(selected_stability) == 1,
           round(selected_stability, 4), "NOT FOUND (NA)"), "\n")
cat("Jaccard threshold (0.40)        :",
    ifelse(length(selected_stability) == 1 && selected_stability >= 0.40,
           "MET", "NOT MET"), "\n\n")

fallback_triggered <- FALSE
if (length(selected_stability) == 1 &&
    !is.na(selected_stability) &&
    selected_stability < 0.5 &&
    nrow(stability_summary) > 1) {

  alt_k <- stability_summary[order(-jaccard, -rand, k)][1, k]
  alt_jaccard <- stability_summary[order(-jaccard, -rand, k)][1, jaccard]
  cat("Fallback condition met. Evaluating stability-optimal k:", alt_k, "\n")
  cat("Jaccard at stability-optimal k:", round(alt_jaccard, 4), "\n")

  if (alt_jaccard > selected_stability) {
    selected_k     <- alt_k
    fallback_triggered <- TRUE
    cat("Fallback ACTIVATED: selected_k changed to", selected_k, "\n\n")
  } else {
    selected_k <- validation_optimal_k_raw
    cat("Fallback NOT activated: stability-optimal k has no better Jaccard.\n")
    cat("Retaining silhouette-optimal k =", selected_k, "\n\n")
  }
} else {
  selected_k <- validation_optimal_k_raw
  cat("Fallback condition not met. Using silhouette-optimal k =",
      selected_k, "\n\n")
}
cat("FINAL selected_k :", selected_k, "\n")
cat("========================================\n")
\end{lstlisting}

\subsection{Change 6-B: Quality Gate Warning Block (Critical)}

\textbf{Location:} Immediately after the k-selection block above
(after the diagnostic logging, before \rcode{selected\_model} is assigned).\\
\textbf{Audit source:} Defect 1 in \rfile{Li\_Clustering\_Audit\_Report.tex}

\textbf{New code to insert:}
\begin{lstlisting}
# --- v2.0: Explicit cluster quality gate ---
selected_silhouette <- validation_table[k == selected_k, silhouette]
selected_jaccard_v2 <- if (selected_k %in% stability_summary$k) {
  stability_summary[k == selected_k, jaccard]
} else {
  NA_real_
}

sil_pass     <- !is.na(selected_silhouette) && selected_silhouette >= 0.30
jaccard_pass <- !is.na(selected_jaccard_v2) && selected_jaccard_v2 >= 0.40

quality_gate <- data.table(
  metric         = c("Silhouette width", "Jaccard stability"),
  observed       = c(round(selected_silhouette, 4),
                     round(selected_jaccard_v2, 4)),
  threshold      = c(0.30, 0.40),
  pass           = c(sil_pass, jaccard_pass),
  interpretation = c(
    ifelse(sil_pass,
           "Acceptable cluster separation",
           "Weak cluster separation -- exploratory status"),
    ifelse(jaccard_pass,
           "Acceptable bootstrap stability",
           "Poor bootstrap stability -- exploratory status")
  )
)

fwrite(quality_gate, file.path(output_dir, "phase6_quality_gate.csv"))
knitr::kable(quality_gate, caption = "Phase 6 cluster quality gate")

if (!sil_pass || !jaccard_pass) {
  warning(paste0(
    "\n\n*** CLUSTER QUALITY WARNING (v2.0) ***\n",
    "  Silhouette = ", round(selected_silhouette, 4),
    " (threshold >= 0.30): ", ifelse(sil_pass, "PASS", "FAIL"), "\n",
    "  Jaccard    = ", round(selected_jaccard_v2, 4),
    " (threshold >= 0.40): ", ifelse(jaccard_pass, "PASS", "FAIL"), "\n",
    "  The selected k = ", selected_k,
    " scenario clustering solution is EXPLORATORY.\n",
    "  Do not present discovered clusters as stable accident archetypes.\n",
    "  Present findings as candidate patterns for hypothesis generation.\n",
    "***********************************************\n"
  ))
}
\end{lstlisting}

\subsection{Change 6-C: Jaccard Confidence Interval Under Consensus Profile (Medium)}

\textbf{Location:} Phase 6 stability summary block (lines $\approx$1240--1264).\\
\textbf{Audit source:} Improvement IP-3

\textbf{Problem:} Under \rcode{local\_safe} (one repeat), Jaccard is a single
point estimate with no variance information. Under \rcode{consensus\_16gb}
(3 repeats), multiple Jaccard values are available and should be summarized
with mean and standard deviation.

\textbf{Change:} After building \rcode{stability\_summary}, add:
\begin{lstlisting}
# v2.0: If multiple stability repeats are available, compute Jaccard SD
if (!is.null(stability_repeat_results) && nrow(stability_repeat_results) > 0) {
  jaccard_ci <- stability_repeat_results[, .(
    jaccard_mean = round(mean(jaccard, na.rm = TRUE), 4),
    jaccard_sd   = round(sd(jaccard, na.rm = TRUE), 4),
    jaccard_min  = round(min(jaccard, na.rm = TRUE), 4),
    jaccard_max  = round(max(jaccard, na.rm = TRUE), 4),
    n_repeats    = .N
  ), by = k]

  fwrite(jaccard_ci, file.path(output_dir, "phase6_jaccard_ci.csv"))
  knitr::kable(jaccard_ci,
               caption = "Phase 6 Jaccard stability with confidence range")
} else {
  cat("Single bootstrap repeat --- Jaccard CI not available under local_safe.\n")
}
\end{lstlisting}

\subsection{Change 6-D: Reduced-Feature Validation (High)}

\textbf{Location:} After the primary validation block, still in Phase 6.

\textbf{Purpose:} Run \rcode{validation\_kproto()} on the reduced-feature
candidate models to compare silhouette scores with and without
\rcode{Pressure(in)} and \rcode{Visibility(mi)}.

\textbf{New chunk --- name it \rcode{phase-6-reduced-validation}:}
\begin{lstlisting}
# v2.0: Validate reduced-feature k-prototypes candidates
set.seed(42)
validation_reduced_raw <- clustMixType::validation_kproto(
  method  = "silhouette",
  k       = candidate_k_values,
  x       = df_scenario_reduced,
  nstart  = 5,
  type    = "gower",
  verbose = FALSE
)

validation_reduced_table <- as.data.table(validation_reduced_raw$indices)
setnames(validation_reduced_table,
         c("clusters", "silhouette"),
         c("k", "silhouette"))
setorder(validation_reduced_table, -silhouette)
validation_reduced_table[, feature_set := "reduced"]

# Combine with full-feature silhouette table for comparison
validation_full_labeled <- copy(validation_table)[
  , .(k, silhouette)
][, feature_set := "full"]

feature_silhouette_comparison <- rbind(
  validation_full_labeled, validation_reduced_table[, .(k, silhouette, feature_set)]
)
setorder(feature_silhouette_comparison, feature_set, -silhouette)

fwrite(feature_silhouette_comparison,
       file.path(output_dir, "phase6_feature_silhouette_comparison.csv"))

knitr::kable(feature_silhouette_comparison,
             caption = "Phase 6: Silhouette comparison -- full vs reduced feature set")
\end{lstlisting}

% ============================================================
\section{Phase 7: DBSCAN Additions}
% ============================================================

\subsection{Change 7-A: Print Cluster Centroid Coordinates (Medium)}

\textbf{Location:} After the DBSCAN execution block and before corridor
label assignment (approximately after line 1564).\\
\textbf{Audit source:} Improvement IP-2

\textbf{Purpose:} Provide centroid latitude and longitude for each spatial
cluster so corridor assignments can be verified manually in a mapping tool.

\textbf{New code to insert:}
\begin{lstlisting}
# v2.0: Compute and export spatial cluster centroids for corridor verification
cluster_centroids <- df_spatial[spatial_cluster > 0, .(
  centroid_x_km  = round(mean(x_km), 4),
  centroid_y_km  = round(mean(y_km), 4),
  centroid_lat   = round(
    40.8 + mean(y_km) / 110.574, 5
  ),
  centroid_lng   = round(
    -73.2 + mean(x_km) / (111.32 * cos(40.8 * pi / 180)), 5
  ),
  n_accidents = .N
), by = spatial_cluster]

setorder(cluster_centroids, spatial_cluster)

fwrite(cluster_centroids,
       file.path(output_dir, "phase7_cluster_centroids.csv"))

knitr::kable(cluster_centroids,
  caption = paste(
    "Phase 7 spatial cluster centroids.",
    "Verify each centroid against the named corridor.",
    "Reference: lat0 = 40.8 N, lng0 = 73.2 W."
  )
)
\end{lstlisting}

\textbf{Verification instruction:} After the document renders, open
\rfile{phase7\_cluster\_centroids.csv} and paste each
\rcode{centroid\_lat} / \rcode{centroid\_lng} pair into a map. Confirm
that the point falls on or within $\approx$1 km of the named highway.
If any assignment is wrong, update the latitude band thresholds at
approximately lines 1584--1591 of the Rmd and re-render.

\subsection{Change 7-B: eps Sensitivity Check (Medium)}

\textbf{Location:} After the primary DBSCAN run (after line $\approx$1564).\\
\textbf{Audit source:} Improvement IP-6

\textbf{New chunk --- name it \rcode{phase-7-eps-sensitivity}:}
\begin{lstlisting}
# v2.0: Test eps sensitivity at +/- 0.1 km from selected value
eps_sensitivity_results <- lapply(
  c(selected_eps - 0.1, selected_eps, selected_eps + 0.1),
  function(eps_test) {
    res <- dbscan::dbscan(
      df_spatial[, .(x_km, y_km)],
      eps     = eps_test,
      minPts  = selected_minPts
    )
    n_clusters <- length(unique(res$cluster[res$cluster > 0]))
    noise_pct  <- round(100 * mean(res$cluster == 0), 2)
    max_share  <- round(100 * max(table(res$cluster[res$cluster > 0])) /
                          sum(res$cluster > 0), 2)
    data.table(
      eps          = eps_test,
      n_clusters   = n_clusters,
      noise_pct    = noise_pct,
      max_cluster_share_pct = max_share,
      within_srs_targets = (n_clusters >= 3 & n_clusters <= 15 &
                              noise_pct >= 5 & noise_pct <= 30 &
                              max_share <= 60)
    )
  }
)
eps_sensitivity_table <- rbindlist(eps_sensitivity_results)

fwrite(eps_sensitivity_table,
       file.path(output_dir, "phase7_eps_sensitivity.csv"))

knitr::kable(eps_sensitivity_table,
             caption = "Phase 7 DBSCAN eps sensitivity check (+/-0.1 km)")
\end{lstlisting}

% ============================================================
\section{Phase 8: Interpretation Corrections}
% ============================================================

\subsection{Change 8-A: Fix Cluster 2 Naming (Medium)}

\textbf{Location:} Lines 1894--1912 (chunk \rcode{phase-8-final-names}).\\
\textbf{Audit source:} Improvement IP-4

\textbf{Problem:} The naming hierarchy never matches Cluster 2 with a
meaningful condition, so it receives the generic fallback name
``Scenario Cluster 2 --- Cloudy.'' Cluster 2 has 3{,}857 rows and Sev3/4 of
17.27\%, making it the second-largest risk group. It deserves a real name.

\textbf{Diagnosis first:} Before changing the naming code, add a diagnostic
print block to see Cluster 2's actual profile values. Insert after the
\rcode{scenario\_weather\_time\_profile} is built:

\begin{lstlisting}
# v2.0: Diagnostic -- inspect all profile values to guide naming
cat("=== Cluster Profile Values for Naming Diagnostics ===\n")
print(scenario_weather_time_profile[, .(
  scenario_cluster, dominant_hour_band, dominant_weather_group,
  night_pct, weekend_pct, junction_pct, traffic_signal_pct,
  crossing_pct
)])
cat("=====================================================\n")
\end{lstlisting}

\textbf{Updated naming hierarchy} (replace lines 1894--1912):
\begin{lstlisting}
scenario_final_names[, final_name := dplyr::case_when(
  # Tier 1: Extreme weather
  dominant_weather_group == "Snow_Ice" ~
    "Winter Adverse --- Snow and Ice",

  # Tier 2: Night with hazardous conditions
  night_pct >= 50 &
    dominant_weather_group %in% c("Fog_Haze", "Rain", "Thunderstorm") ~
    "Night Low-Visibility --- Fog and Rain",

  # Tier 3: Night / early morning (any weather)
  night_pct >= 50 & dominant_hour_band == "Early_AM" ~
    paste("Night / Early AM ---",
          gsub("_", " ", dominant_weather_group)),

  # Tier 4: Weekend surface road
  weekend_pct >= 40 & traffic_signal_pct >= 8 ~
    "Weekend Surface Road --- Signal Heavy",

  # Tier 5: Rush corridor
  dominant_hour_band %in% c("AM_Rush", "PM_Rush") &
    junction_pct >= 10 ~
    paste(gsub("_", " ", dominant_weather_group),
          gsub("_", " ", dominant_hour_band),
          "--- Junction Heavy"),

  # Tier 6 (v2.0 NEW): Cloudy midday surface road
  # Captures clusters dominated by Cloudy weather in off-peak daytime hours
  dominant_weather_group == "Cloudy" &
    dominant_hour_band %in% c("Midday", "Evening") ~
    paste("Cloudy Off-Peak ---",
          gsub("_", " ", dominant_hour_band)),

  # Tier 7 (v2.0 NEW): Clear daytime non-rush
  dominant_weather_group == "Clear" &
    dominant_hour_band %in% c("Midday", "Evening") ~
    paste("Clear Off-Peak ---",
          gsub("_", " ", dominant_hour_band)),

  # Fallback: still generic but more informative
  TRUE ~
    paste("Mixed Conditions ---",
          gsub("_", " ", dominant_weather_group),
          gsub("_", " ", dominant_hour_band))
)]
\end{lstlisting}

\textbf{Note:} Run the diagnostic print first. The Tier 6 and Tier 7
conditions use \rcode{dominant\_hour\_band} values from the actual profile.
If Cluster 2's dominant hour band is neither Midday nor Evening, adjust
the condition to match the observed value.

\subsection{Change 8-B: Add Explicit Severity 4 Rate Column (Medium)}

\textbf{Location:} Phase 8 severity profiling block (approximately
lines 1769--1790).\\
\textbf{Audit source:} Improvement IP-5

\textbf{Change:} After building \rcode{scenario\_severity\_profile},
add a baseline comparison row and export the expanded table:

\begin{lstlisting}
# v2.0: Add sev4_rate column and baseline comparison row
dataset_sev4_rate <- round(100 * mean(df_clustered$Severity == 4), 2)
dataset_sev34_rate <- round(
  100 * mean(df_clustered$Severity %in% c(3, 4)), 2
)

scenario_severity_profile[, sev4_rate := round(
  100 * sev4_count / n_accidents, 2
)]
scenario_severity_profile[, sev4_vs_baseline := round(
  sev4_rate / dataset_sev4_rate, 2
)]

# Baseline row for easy comparison
baseline_row <- data.table(
  scenario_cluster = 0L,
  n_accidents      = nrow(df_clustered),
  sev1_pct = round(100 * mean(df_clustered$Severity == 1), 2),
  sev2_pct = round(100 * mean(df_clustered$Severity == 2), 2),
  sev3_pct = round(100 * mean(df_clustered$Severity == 3), 2),
  sev4_pct = round(100 * mean(df_clustered$Severity == 4), 2),
  sev34_pct = dataset_sev34_rate,
  sev4_count = sum(df_clustered$Severity == 4),
  high_risk_archetype = FALSE,
  sev4_rate    = dataset_sev4_rate,
  sev4_vs_baseline = 1.00
)

severity_profile_v2 <- rbind(
  baseline_row,
  scenario_severity_profile,
  fill = TRUE
)

fwrite(severity_profile_v2,
       file.path(output_dir, "phase8_severity_profile_v2.csv"))

knitr::kable(severity_profile_v2,
  caption = paste(
    "Phase 8 severity profile with Sev4 rate per cluster.",
    "Row 0 = dataset baseline.",
    "sev4_vs_baseline = ratio of cluster Sev4 rate to baseline."
  )
)
\end{lstlisting}

% ============================================================
\section{Phase 9: Export and Visualization Corrections}
% ============================================================

\subsection{Change 9-A: Add Warning Annotation to Stability Plot (Medium)}

\textbf{Location:} Phase 9 stability plot chunk (approximately
lines 2101--2158).\\
\textbf{Audit source:} Defect 1 (visualization component)

\textbf{Problem:} The stability bar chart shows all bars in the red zone
but carries no label or subtitle informing the reader why this is
significant. A reader encountering the plot in the final report without
reading the surrounding text would not understand the implication.

\textbf{Change:} Find the \rcode{ggplot()} call that builds the stability
plot. Modify the \rcode{labs()} call to include an explicit subtitle:

\begin{lstlisting}
# v2.0: Add quality-status subtitle to stability plot
stability_quality_label <- if (!sil_pass || !jaccard_pass) {
  paste0(
    "WARNING: Results are EXPLORATORY. ",
    "Jaccard = ", round(selected_jaccard_v2, 3),
    " (threshold 0.40). ",
    "Silhouette = ", round(selected_silhouette, 3),
    " (threshold 0.30)."
  )
} else {
  "Cluster quality meets acceptable thresholds."
}

# Within the existing ggplot call, update labs():
# Replace the existing labs() with:
# labs(
#   title    = "Phase 6 Cluster Stability --- Rand and Jaccard by k",
#   subtitle = stability_quality_label,
#   x = "k (number of clusters)",
#   y = "Index value",
#   fill = "Metric"
# )
\end{lstlisting}

\textbf{Implementation note:} \rcode{sil\_pass} and \rcode{jaccard\_pass}
are defined in Change 6-B. The Phase 9 chunk runs after Phase 6 in the
same R session, so these variables will be available.

\subsection{Change 9-B: Feature Sensitivity Comparison Export (High)}

\textbf{Location:} Phase 9 export block, after existing table exports.

\textbf{New code:}
\begin{lstlisting}
# v2.0: Export feature sensitivity comparison
if (file.exists(file.path(output_dir, "phase6_feature_silhouette_comparison.csv"))) {
  feature_sil_compare <- fread(
    file.path(output_dir, "phase6_feature_silhouette_comparison.csv")
  )

  p_feature_compare <- ggplot(
    feature_sil_compare,
    aes(x = factor(k), y = silhouette, fill = feature_set)
  ) +
    geom_col(position = "dodge", width = 0.6) +
    geom_hline(yintercept = 0.30, linetype = "dashed",
               color = "red", linewidth = 0.8) +
    annotate("text", x = 0.7, y = 0.31, hjust = 0, size = 3,
             color = "red", label = "Threshold 0.30") +
    scale_fill_manual(
      values = c("full" = "#4472C4", "reduced" = "#ED7D31"),
      labels = c("Full (16 features)", "Reduced (14 features)")
    ) +
    labs(
      title    = "Phase 6 Silhouette: Full vs Reduced Feature Set",
      subtitle = "Reduced set drops Pressure(in) and Visibility(mi)",
      x = "k", y = "Silhouette width", fill = "Feature set"
    ) +
    theme_bw(base_size = 11)

  ggsave(
    file.path(output_dir, "figures", "phase9_feature_sensitivity.png"),
    plot   = p_feature_compare,
    width  = 7, height = 4, dpi = 300
  )
}
\end{lstlisting}

% ============================================================
\section{Phase 10: Limitations Table Correction}
% ============================================================

\subsection{Change 10-A: Add Cluster Quality Failure to Limitations (Critical)}

\textbf{Location:} Phase 10 limitations table block (approximately
lines 2756--2790).\\
\textbf{Audit source:} Defect 1 (report transparency component)

\textbf{Problem:} The existing limitations table lists six items including
``sampled validation,'' duration capping, and precipitation exclusion.
It does not contain any row stating that the scenario cluster quality
metrics are below standard thresholds or that the scenario results are
therefore exploratory. A reader of only the report PDF has no way to
know that Sanity Check 4 failed.

\textbf{Change:} Locate the \rcode{phase10\_limitations\_table} construction.
Add the following rows to the data.table:

\begin{lstlisting}
# v2.0: Add cluster quality limitation rows to the limitations table
# These rows go FIRST so they are immediately visible in the rendered table.
cluster_quality_limitations <- data.table(
  limitation = c(
    "Scenario cluster quality below standard thresholds",
    "Scenario clustering presented as exploratory"
  ),
  detail = c(
    paste0(
      "Silhouette width = ",
      round(selected_silhouette, 3),
      " (published threshold >= 0.30, Kaufman and Rousseeuw 1990). ",
      "Jaccard stability index = ",
      round(selected_jaccard_v2, 3),
      " (threshold >= 0.40 for exploratory, >= 0.60 for acceptable, ",
      "Hennig 2007). Both metrics fall below acceptable thresholds."
    ),
    paste0(
      "Because silhouette and Jaccard are both below threshold, ",
      "the k = ", selected_k, " scenario clustering solution should be ",
      "interpreted as a set of candidate patterns for hypothesis generation, ",
      "not as stable or confirmed accident archetypes. ",
      "The DBSCAN spatial clustering results are separately validated and ",
      "remain defensible."
    )
  ),
  impact = c(
    "Cluster boundaries are weak; different samples may produce different partitions.",
    "Severity 3/4 elevation in Clusters 4 and 5 is a plausible signal but requires replication."
  ),
  mitigation = c(
    "Rerun with full_srs profile on a high-memory machine to obtain full-data silhouette estimates.",
    "Present findings with explicit exploratory caveats in all report sections discussing scenario clusters."
  )
)

# Prepend these rows before the existing limitations
phase10_limitations_table <- rbind(
  cluster_quality_limitations,
  phase10_limitations_table,
  fill = TRUE
)
\end{lstlisting}

% ============================================================
\section{Verification Checklist After All Changes}
% ============================================================

After implementing all changes in \rfile{li\_clustering\_v2.Rmd} and
successfully knitting the document, verify the following:

\begin{longtable}{|p{0.5cm}|p{4.5cm}|p{8.5cm}|}
\hline
\textbf{\#} & \textbf{Check} & \textbf{Expected outcome} \\
\hline
\endfirsthead
\hline
\textbf{\#} & \textbf{Check} & \textbf{Expected outcome} \\
\hline
\endhead
1 & Profile is consensus\_16gb & Line 135 of v2 reads \rcode{run\_profile <- "consensus\_16gb"} \\
\hline
2 & Quality gate CSV exists & \rcode{phase6\_quality\_gate.csv} in output directory \\
\hline
3 & Warning fires if Jaccard < 0.40 & R console shows the WARNING block during render \\
\hline
4 & k-selection trace printed & Console output shows ``=== Phase 6 k-Selection Diagnostics ==='' \\
\hline
5 & Jaccard CI table renders & \rcode{phase6\_jaccard\_ci.csv} populated (if consensus profile) \\
\hline
6 & Reduced feature matrix built & \rcode{df\_scenario\_reduced} has 14 columns, zero NAs \\
\hline
7 & Feature sensitivity CSV exists & \rcode{phase6\_feature\_silhouette\_comparison.csv} \\
\hline
8 & Cluster 2 has non-generic name & \rcode{phase8\_scenario\_final\_names.csv} contains no row matching ``Scenario Cluster 2 ---'' \\
\hline
9 & Severity 4 column present & \rcode{phase8\_severity\_profile\_v2.csv} has \rcode{sev4\_rate} and \rcode{sev4\_vs\_baseline} columns; row 0 is baseline \\
\hline
10 & Centroid CSV exists & \rcode{phase7\_cluster\_centroids.csv} with lat/lng columns \\
\hline
11 & eps sensitivity CSV exists & \rcode{phase7\_eps\_sensitivity.csv} with 3 rows \\
\hline
12 & Stability plot has subtitle & Rendered PDF stability chart shows the quality warning subtitle \\
\hline
13 & Limitations table has quality rows & First two rows of \rcode{phase10\_limitations\_table} address silhouette and Jaccard failure \\
\hline
14 & Sanity audit check 4 & Still expected to show FAIL unless consensus profile improves Jaccard; result must be accurately reported \\
\hline
15 & Binary flags use native levels & No \rcode{labels = c("FALSE","TRUE")} pattern anywhere in v2 file \\
\hline
\end{longtable}

% ============================================================
\section{Decision Tree After Seeing Consensus Profile Results}
% ============================================================

Once \rfile{li\_clustering\_v2.Rmd} has been knitted under
\rcode{consensus\_16gb}, the following decision tree determines next steps:

\subsection{If consensus Jaccard $\geq 0.40$}

The \rcode{local\_safe} Jaccard of 0.1989 was a sampling artifact.
The full consensus result is acceptable.

\begin{itemize}
\item Remove the exploratory disclaimers added in Changes 6-B, 9-A, 10-A.
\item Update the limitations table to say: ``Initial sampled validation
  under-estimated stability; consensus validation confirms acceptable
  cluster quality (Jaccard $\geq 0.40$).''
\item Present the scenario clustering results as meaningful, not merely
  exploratory.
\end{itemize}

\subsection{If consensus Jaccard remains $< 0.40$}

The cluster instability is real and not a sampling artifact.

\begin{itemize}
\item Keep all exploratory disclaimers exactly as written.
\item Check the feature sensitivity comparison
  (\rcode{phase6\_feature\_silhouette\_comparison.csv}). If the reduced
  feature set (without \rcode{Pressure(in)} and \rcode{Visibility(mi)})
  shows meaningfully higher silhouette (improvement $> 0.03$), switch the
  primary scenario pipeline to use \rcode{df\_scenario\_reduced} as input.
\item If neither full nor reduced feature sets reach Jaccard $\geq 0.40$,
  the data does not strongly support discrete accident scenario archetypes.
  The report must say so plainly and pivot emphasis to the DBSCAN spatial
  results, which are valid and defensible.
\end{itemize}

\subsection{If Cluster 2 still gets a generic name after Change 8-A}

Run the diagnostic print block from Change 8-A. Inspect the actual
\rcode{dominant\_hour\_band} value for Cluster 2. Add a Tier 6 condition
that matches that exact value. Re-render.

\subsection{If Cluster 2 name duplicates another cluster after naming fix}

The deduplication logic at approximately line 1921--1924 appends
\rcode{(S\#)} suffixes automatically. Verify that the suffix appears in
the output CSV and update the corresponding Phase 8 narrative in the
final report.

% ============================================================
\section{References}
% ============================================================

\begin{thebibliography}{9}

\bibitem{kaufman1990}
Leonard Kaufman and Peter J.\ Rousseeuw.
\textit{Finding Groups in Data: An Introduction to Cluster Analysis}.
Wiley, 1990.

\bibitem{hennig2007}
Christian Hennig.
\textit{Cluster-wise assessment of cluster stability}.
Computational Statistics and Data Analysis, 52(1):258--271, 2007.

\bibitem{szepannek2018}
Gero Szepannek.
\textit{clustMixType: User-Friendly Clustering of Mixed-Type Data in R}.
The R Journal, 2018.
\url{https://journal.r-project.org/articles/RJ-2018-048/}

\bibitem{hahsler2019}
Michael Hahsler, Matthew Piekenbrock, and Derek Doran.
\textit{dbscan: Fast Density-Based Clustering with R}.
Journal of Statistical Software, 91(1), 2019.
\url{https://www.jstatsoft.org/v91/i01/}

\end{thebibliography}

\end{document}
